% =====================================================================
% sections/01-abstract.tex – BORRADOR (estructura BMC: Antecedentes,
% Resultados, Conclusiones; revisar el límite de palabras en JOURNAL.md)
% Los [ ] se llenan con los resultados del experimento de k-mers.
% =====================================================================

\textbf{Antecedentes.} Los elementos transponibles (TEs) se cuentan
entre los componentes más abundantes de los genomas eucariotas, y su
estudio combina datos de naturaleza muy distinta: secuencias, una
clasificación en tres niveles, la taxonomía de cada especie y perfiles
de \textit{k}-mers de hasta 5\,460 valores por secuencia. Recursos como
PanTEon distribuyen esta información en archivos planos unidos solo por
el nombre de la especie, sin consultas ni relaciones explícitas. El
objetivo de este trabajo fue diseñar, implementar y evaluar una base de
datos documental en MongoDB que integre esas fuentes.

\textbf{Resultados.} Un \textit{pipeline} reproducible en Python leyó
y validó 239\,946 secuencias de 2\,790 especies de Animalia, Plantae y
Fungi, corrigió seis problemas de calidad del FASTA (entre ellos
secuencias en minúsculas, nombres de especie que no cruzaban con el CSV
y 22,9\,\% de secuencias con bases indeterminadas) sin descartar
ningún registro, y calculó los perfiles de \textit{k}-mers de $k=1$ a
6 (1\,439\,676 documentos). El modelo usa tres colecciones relacionadas
por referencias, seis índices justificados por consultas concretas y
validación \texttt{\$jsonSchema} en todas las colecciones. Las
consultas mostraron una base dominada por retrotransposones LTR
(49,7\,\%), con perfiles distintos por reino (LTR en plantas, LINE en
animales), y revelaron que 45 entradas registradas como especie son en
realidad taxones superiores. Al comparar dos estrategias de
almacenamiento de los \textit{k}-mers, la colección aparte tardó
[TIEMPO\_1]\,s en cargarse y ocupó [ESPACIO\_1]\,MB, frente a
[TIEMPO\_2]\,s y [ESPACIO\_2]\,MB de la estrategia embebida.

\textbf{Conclusiones.} Un modelo documental con referencias permite
integrar secuencias, taxonomía y perfiles de \textit{k}-mers sin
repetir información y responder las consultas con agregaciones. La
elección entre guardar los \textit{k}-mers aparte o embebidos depende
del tipo de uso: [CONCLUSIÓN DEL EXPERIMENTO, p. ej. la colección aparte
conviene para consultar un $k$ concreto y la embebida para recuperar el
perfil completo de una secuencia].

\medskip
\noindent\textbf{Palabras clave:} elementos transponibles, MongoDB,
bases de datos documentales, \textit{k}-mers, bioinformática.
